% Part: proof-theory
% Chapter: normalization
% Section: normalization-thm

\documentclass[../../../include/open-logic-section]{subfiles}

\begin{document}

\olfileid{pt}{nor}{thm}

\olsection{Teorema Normalisasi}

!!^a{proof} ana ing wujud \emph{normal} yèn ora nduwèni segmen cut.
Cetha yèn iki kelakon yèn lan mung yèn ora ana konversi permutasi
utawa konversi reduksi kang bisa ditrapaké marang bukti mau. Kanggo
\emph{normalisasi} !!a{proof} tegesé ngowahi bukti mau dadi bukti
normal kanthi ngetrapaké konversi permutasi lan reduksi nganti ora
ana cut kang kari. Pratelan yèn iki tansah bisa ditindakaké diarani
\emph{teorema normalisasi}.

\begin{thm}\ollabel{thm:normalization} Saben \Log{N1i}-!!{proof}
$\delta$ bisa dinormalisasi, yaiku bisa diowahi dadi
!!a{proof}~$\delta'$ tanpa cut.
\end{thm}

\begin{proof}
  Buktiné nganggo induksi marang pangkat cut lan dawa cut
  saka~$\delta$. Hipotesis induksi nyatakaké yèn saben !!{proof}
  kang pangkat cuté luwih cilik, utawa kang pangkat cuté padha
  nanging dawa cuté luwih cilik, bisa dinormalisasi.

  Yèn dawa cuté~$0$, ora ana cut lan $\delta$ wis normal.
  Mula anggepen pangkat cut lan dawa cut saka~$\delta$
  padha-padha~$> 0$. Pilih segmen cut maksimal kang paling
  dhuwur lan paling tengen. Yèn dawa segmen cut mau $>1$,
  trapna konversi permutasi. Yèn dawa segmen mau~$1$,
  trapna konversi reduksi kang cocog. Iki ngasilaké !!{proof}
  anyar kang pangkat cuté padha nanging dawa cuté luwih cilik,
  utawa pangkat cuté luwih cilik; mula hipotesis induksi bisa
  ditrapaké.
\end{proof}

Siasat ing bukti ndhuwur (tansah ngrampungaké cut maksimal kang
paling dhuwur lan paling tengen) menehi urutan tartamtu kanggo
ngetrapaké konversi marang !!a{proof} supaya ngasilaké !!{proof}
normal. Nanging, nggumunaké yèn urutané pancèn ora dadi prakara.
Ngetrapaké konversi permutasi lan reduksi miturut urutan apa waé
pungkasané ngasilaké bukti normal.

Definisi segmen lan cut gampang dirumusaké kanggo \Log{N2i}, lan
konversi permutasi lan reduksi kanggo \Log{N1i} langsung bisa
diterjemahaké menyang \Log{N2i}. Mula teorema normalisasi uga
lumaku kanggo \Log{N2i}.

\begin{thm}\ollabel{thm:normalization-N2i} Saben \Log{N2i}-!!{proof}
$\delta$ bisa dinormalisasi, yaiku bisa diowahi dadi
!!a{proof}~$\delta'$ tanpa cut.
\end{thm}

\end{document}
